\appendix

% =============================================================================
\section{Syntax}
\label{sec:syntax}
% =============================================================================
The Mol-Metal hybrid represents both the symbolic $\lambda$-layer and the
chemistry layer in a single typed term algebra. We follow the nine-layer
catalogue of \texttt{molmetal/reports/lambda\_layer\_metrics.md}, restricted
to the four layers that participate in rewriting: atoms, bonds, applications
and combinators. Geometric priors and empirical scoremix layers are
observational only and do not introduce new syntactic forms.

\subsection{Atoms}
An \emph{atom} $a$ carries a chemical symbol $s \in \Sigma$ (e.g.\ \texttt{C, N, Pt}),
a typed charge $q \in \mathbb{Z}$ and an electron-domain count
$\rho \in \{0,1,2,\dots\}$ (the number of free valences plus lone pairs that
the atom exposes at rewrite time):
\[
  a \;::=\; \mathsf{Atom}(s,\,q,\,\rho).
\]
The domain $\rho$ is the \emph{site arity} used by the FreeSiteLedger in
\texttt{bonds/application.py:FreeSiteLedger}. We require $0 \le \rho \le 6$
for organic atoms and $0 \le \rho \le 8$ for transition metals; Pt(II) and
Pd(II) admit $\rho=4$ and Pt(IV) admits $\rho=6$.

\subsection{Combinators}
The pure-$\lambda$ fragment uses three core nodes:
\begin{align*}
  c \;::=\;& \mathsf{Var}(x,\,\rho)       && \text{typed variable with site count} \\
           \mid\;& \mathsf{Abs}(x{:}\tau,\,t) && \text{abstraction over a typed binder} \\
           \mid\;& \mathsf{Comb}(k;\,\vec a)  && \text{primitive combinator with arguments.}
\end{align*}
The combinator tag $k$ ranges over a fixed symbol set
$\{ \mathsf{CuAAC},\ \mathsf{SPAAC},\ \mathsf{ThiolEne},\ \mathsf{Suzuki},\
\mathsf{AmideCoupling},\ \mathsf{Dative},\ \mathsf{Geometry}_{\mathrm{sp}},\dots\}$
that enumerates every legal rewriting head.

\subsection{Bonds and applications}
A \emph{bond} $b$ is one of six kinds, mirroring RDKit's bond order plus the
dative variant we need for coordination chemistry:
\[
  b \;::=\; \mathsf{single} \mid \mathsf{double} \mid \mathsf{triple}
          \mid \mathsf{aromatic} \mid \mathsf{dative} \mid \mathsf{hydrogen}.
\]
An \emph{application} is the spine $t \, t$ where the head $t$ is a
combinator and the spine is a sequence of atoms or partial molecules.
A \emph{term} $u$ wraps a term $t$ with a typing context $\Gamma$:
\[
  u \;::=\; \mathsf{Term}(t,\,\Gamma).
\]
A \emph{rule} $r$ is a quadruple
$r = (\ell,\,\vec p,\,\vec e,\,\kappa)$ where $\ell$ is the head combinator,
$\vec p$ the preconditions, $\vec e$ the expected effects (atoms consumed and
produced) and $\kappa$ the geometric shape predicate.

\subsection{Types}
\[
  \tau \;::=\; \mathsf{AtomType}(s) \mid \mathsf{BondType}(b)
            \mid \tau \to \tau.
\]
The two base kinds classify the chemistry fragment; the arrow is reserved
for the pure-$\lambda$ fragment.

\subsection{Predicates}
Rewrite rules are guarded by four predicates, evaluated by the
\texttt{well\_formedness.py} layer before any contraction is attempted:
\begin{align*}
  p \;::=\;& \mathsf{HasArity}(a,\,n)        && a \text{ exposes at least } n \text{ sites} \\
           \mid\;& \mathsf{Saturated}(a)        && \text{all sites of } a \text{ are consumed} \\
           \mid\;& \mathsf{Compatible}(a_1,a_2) && \text{electronic structure permits bonding} \\
           \mid\;& \mathsf{Geometry}(a,\tau)    && \text{coordination geometry satisfied.}
\end{align*}
The $\mathsf{Geometry}$ predicate is enforced by
\texttt{priors/metal\_geometry.py:SquarePlanarPtII} for Pt(II)/Pd(II)
($\rho = 4$, $90^\circ$), by $\mathsf{Octahedral}$ for Ru(II)/Ir(III)
($\rho = 6$, $90^\circ$) and by $\mathsf{Tetrahedral}$ for Zn(II).

\begin{example}[\textbf{Azide-alkyne cycloaddition (CuAAC)}]
The rule $r_{\mathsf{CuAAC}}$ has
$\ell = \mathsf{CuAAC}$,
$\vec p = (\mathsf{HasArity}(\text{azide},2),\,
          \mathsf{HasArity}(\text{alkyne},2),\,
          \mathsf{Compatible},\,
          \mathsf{Geometry}(\text{Cu},\mathsf{Tetrahedral}))$,
$\vec e = (\text{loss of 2 sites on azide},\, \text{loss of 2 sites on alkyne},\,
           \text{formation of 5-membered triazole})$
and $\kappa = 1.0$ (no geometric slack).
\end{example}

\subsection{Typing contexts}
$\Gamma$ is a finite map from variables $x$ to types $\tau$ subject to the
usual weakening and exchange properties. We assume a simple type system
$\tau \to \tau$ with two base kinds; this is the simply-typed $\lambda$-calculus
of Church extended with $\mathsf{AtomType}$ and $\mathsf{BondType}$ atoms.

% =============================================================================
\section{Reduction rules}
\label{sec:reduction}
% =============================================================================
We define the \emph{multi-step} reduction relation $\stepsto$ as the
compatible closure of four primitive rewrite heads: $\beta$ (pure lambda),
$\chi$ (click chemistry), $\delta$ (dative application) and $\gamma$
(geometry predicate resolution).

\subsection{Standard $\beta$-reduction}
The pure fragment inherits the classical rule of the $\lambda$-calculus:
\begin{equation*}
  (\lambda x{:}\tau.\, M)\; N \;\stepstoBeta\; \subst{x}{N}{M}
  \qquad\text{(capture-avoiding substitution)}.
\end{equation*}
Reduction is performed modulo $\alpha$-equivalence so that the binders of
$\mathsf{Abs}$ are consistently renamed to avoid clashes
(\texttt{lam\_chem/ast.py:\_alpha\_rename}).

\subsection{Click chemistry rules}
The five click reactions enumerated in \texttt{lam\_chem/rules.py:list\_reactions}
are first-class rewrite rules that consume site pairs and emit a new bond:
\begin{align*}
  \mathsf{CuAAC}(\text{azide},\text{alkyne})
    &\;\stepstoChem\; \text{1,2,3-triazole} \\
  \mathsf{SPAAC}(\text{azide},\text{cyclooctyne})
    &\;\stepstoChem\; \text{1,2,3-triazole (strain-promoted)} \\
  \mathsf{ThiolEne}(\text{thiol},\text{alkene})
    &\;\stepstoChem\; \text{thioether} \\
  \mathsf{Suzuki}(\text{aryl-B(OH)$_2$},\text{aryl-X})
    &\;\stepstoChem\; \text{biaryl} \\
  \mathsf{AmideCoupling}(\text{amine},\text{acid})
    &\;\stepstoChem\; \text{amide}.
\end{align*}
Each rule decreases the \emph{joint site count}
$\sum_a \rho(a)$ by exactly $4$ (CuAAC, SPAAC, Suzuki), by $2$
(ThiolEne) or by $2$ (Amide), and never increases it. The CuAAC and SPAAC
rules preserve atom counts; ThiolEne and Amide emit a small byproduct
(\ensuremath{\mathrm{H_2O}} or \ensuremath{\mathrm{HX}}) that is consumed by the global
\texttt{FreeSiteLedger} but is invisible to the $\lambda$-term view.

\subsection{Dative application}
Dative bonds are introduced by the rule
$\mathsf{Dative}(L,M) \to \mathsf{dative}(L,M)$ where $L$ is a Lewis base
($\rho \ge 1$) and $M$ a metal centre with $\rho \ge 1$. The rule decreases
$\rho(L)$ and $\rho(M)$ by one each, contributing to the same joint site
budget as the click rules. This rule is implemented by
\texttt{bonds/application.py:\_orient\_dative}.

\subsection{Geometry predicate resolution}
The $\gamma$ rewrite is not a producer of new bonds but a \emph{failure
eliminator}: when a partial molecule has saturated the geometry predicate
of a metal centre, $\mathsf{Geometry}(a,\tau)$ collapses to $\mathsf{true}$
and the matching combinator is unlocked. Geometry has no effect on the site
budget; it merely gates the $\chi$ and $\delta$ rules.

\subsection{Compatibility and parallel moves}
When two disjoint redexes coexist in a term, both can fire; when they share a
variable, capture-avoiding substitution forces a sequentialisation. We will
show in Section~\ref{sec:confluence} that this dependency is benign.

% =============================================================================
\section{Strong normalization}
\label{sec:sn}
% =============================================================================
The reduction relation $\stepstoAll$ is strongly normalizing on the set of
well-typed, arity-bounded terms. We prove this by exhibiting a
lexicographic measure into a well-founded multiset order.

\subsection{A measure on terms}

\begin{definition}[Site-weighted redex count]
\label{def:measure}
For a term $u = \mathsf{Term}(t,\Gamma)$, define the triple
\[
  \mu(u) \;=\; \bigl(\, n_\beta(u),\; n_{\text{site}}(u),\; n_{\text{chem}}(u) \,\bigr)
\]
where
\begin{itemize}
  \item $n_\beta(u)$ is the number of $(\lambda x.M)N$ redexes in $t$,
  \item $n_{\text{site}}(u)$ is the joint free-site count
        $\sum_a \rho(a)$ summed over every atom reachable from $t$,
  \item $n_{\text{chem}}(u)$ is the number of unmatched functional groups
        (azide, alkyne, thiol, etc.) that no rewrite rule can consume
        without violating a geometry predicate.
\end{itemize}
\end{definition}

\begin{lemma}[Termination of $\beta$]
\label{lem:beta-sn}
The pure $\lambda$-fragment is strongly normalizing under $\stepstoBeta$
within the simply-typed discipline.
\end{lemma}
\begin{proof}
Simply-typed $\lambda$-terms have no fixpoint operator and every $\beta$-step
strictly reduces the size of the derivation tree. By the standard
metatheorem (see \cite[Chapter~6]{Pierce2002}) every closed, well-typed
simply-typed $\lambda$-term is $\sn$.
\end{proof}

\begin{lemma}[Termination of $\chi \cup \delta$]
\label{lem:chem-sn}
Each click-chemistry rule and each dative application strictly decreases
$n_{\text{site}}$ by at least $2$.
\end{lemma}
\begin{proof}
Direct inspection: every $\chi$-rule consumes at least one bond pair from
two distinct atoms; the dative rule consumes one site from the donor and
one from the acceptor. The FreeSiteLedger records these as
\texttt{AtomSite.consume} operations; the post-condition is that
$\rho(a) \ge \rho'(a) + 1$ for at least one participating atom and similarly
for the partner. Hence $n_{\text{site}}(u') \le n_{\text{site}}(u) - 2$.
\end{proof}

\begin{lemma}[Termination of $\gamma$]
\label{lem:gamma-sn}
The $\gamma$ rewrite is locally terminating and globally side-effect-free.
\end{lemma}
\begin{proof}
$\gamma$ is a pure predicate: it collapses a guard
$\mathsf{Geometry}(a,\tau)$ to $\mathsf{true}$. The collapse has no site
effect and does not produce new $\beta$- or $\chi$-redexes. Hence
$\mu(\gamma(u)) \le_{\mathsf{lex}} \mu(u)$ and $\mu$ never increases under
$\gamma$. Because the predicate is decidable and its resolution terminates
(\texttt{priors/metal\_geometry.py:GeometryPenalty} is finite-state on the
bounded geometry classes), $\gamma$ itself has no infinite chain.
\end{proof}

\begin{theorem}[Strong normalization of $\bnf$]
\label{thm:sn}
Every well-typed arity-bounded term is strongly normalizing under
$\stepstoAll = \stepstoBeta \cup \stepstoChem \cup \stepsto_{\delta} \cup
\stepsto_{\gamma}$.
\end{theorem}
\begin{proof}
We lift $\mu$ to the multiset extension over the immediate redex set of $u$:
\[
  M(u) \;=\; \{\, \mu(r) \mid r \text{ is a redex in } u \,\}.
\]
The multiset $M(u)$ is ordered by the Dershowitz--Manna multiset
ordering induced by the lexicographic order on $\mu$. The lexicographic
order is well founded because each component is bounded below by $0$ and
no infinite descending chain can lower the same component twice without
switching components.

By Lemma~\ref{lem:beta-sn}, a $\beta$-step replaces one $(\lambda x.M)N$
redex by the (possibly empty) set of redexes inside $M[x:=N]$; because
$M[x:=N]$ has strictly fewer $\beta$-redexes, $\mu$ decreases in the first
component. By Lemma~\ref{lem:chem-sn}, a $\chi$- or $\delta$-step decreases
$n_{\text{site}}$ by at least $2$, so $\mu$ decreases in the second
component. By Lemma~\ref{lem:gamma-sn}, a $\gamma$-step never increases
$\mu$. No rewrite head ever increases any component.

Since the multiset extension of a well-founded order is itself well
founded (Dershowitz and Manna, 1979; see also \cite{Girard1989}), $M(u)$
admits no infinite descending chain. Hence $\stepstoAll$ is $\sn$.
\end{proof}

\begin{corollary}[Decidability of type-checking]
Type-checking in $\bnf$ is decidable and admits a normalising evaluator.
\end{corollary}
\begin{proof}
By Theorem~\ref{thm:sn} every reduction chain is finite; the normal form
$\mathsf{nf}(u)$ is unique by Theorem~\ref{thm:confluence} below. A
terminating evaluator exists by recursion along $\mu$.
\end{proof}

% =============================================================================
\section{Confluence (Church--Rosser)}
\label{sec:confluence}
% =============================================================================

\begin{theorem}[Church--Rosser property]
\label{thm:confluence}
If $u \stepstoAll^* v_1$ and $u \stepstoAll^* v_2$ then there exists
$w$ such that $v_1 \stepstoAll^* w$ and $v_2 \stepstoAll^* w$.
\end{theorem}

The proof has three layers: the classical Tait--Martin-L\"of argument for
$\beta$, an extension to the disjoint-redex parallel moves of $\chi$ and
$\delta$, and an appeal to the fact that $\gamma$ is observationally
invisible.

\subsection{Local confluence of $\beta$}

\begin{lemma}[$\beta$ has the diamond property]
\label{lem:beta-diamond}
Let $r_1, r_2$ be two distinct $\beta$-redexes in $u$. There exists $w$
with $u \stepstoBeta^* w$ via either order in at most $2$ $\beta$-steps.
\end{lemma}
\begin{proof}
Two cases: $r_1$ and $r_2$ are disjoint (their bound variables are not
shared and they live in separate subterms) or one contains the other. In
the disjoint case each redex fires independently and the two results join
back in two steps. In the nested case the inner redex fires first, after
which the outer one fires (or vice versa by $\alpha$-renaming). This is the
classical parallel-moves argument of Tait and Martin-L\"of; see
\cite[Chapter~3]{Barendregt1984}.
\end{proof}

\begin{lemma}[Local confluence of $\beta$]
\label{lem:beta-lc}
$\stepstoBeta$ is locally confluent.
\end{lemma}
\begin{proof}
Immediate from Lemma~\ref{lem:beta-diamond}: every pair of one-step
divergences from a common ancestor joins back. By Newman's lemma (the
reduction is terminating by Lemma~\ref{lem:beta-sn}) local confluence
implies confluence. Hence $\stepstoBeta$ is $\crprop$.
\end{proof}

\subsection{Parallel moves for $\chi$ and $\delta$}

\begin{lemma}[$\chi \cup \delta$ has the parallel-move property]
\label{lem:chem-diamond}
Let $r_1, r_2$ be $\chi$- or $\delta$-redexes in $u$. If $r_1$ and $r_2$
share an atom, then the rule \texttt{disjoint-rewrites-stall} requires them
to be sequentialised; otherwise they commute.
\end{lemma}
\begin{proof}
Two $\chi$ redexes that share an atom cannot both fire because
$\rho(a)$ cannot decrease twice; the second must wait. Two $\chi$ redexes
that touch disjoint atoms fire independently in either order: applying
$\mathsf{CuAAC}$ to (azide$_1$, alkyne$_1$) and $\mathsf{ThiolEne}$ to
(thiol$_2$, alkene$_2$) leaves the atoms unchanged. The result is
associative: $\mathsf{CuAAC}\bigl(\mathsf{ThiolEne}(u)\bigr)
       = \mathsf{ThiolEne}\bigl(\mathsf{CuAAC}(u)\bigr)$
because both reduce the joint site count by $2$ and neither depends on
the other's byproduct. The same argument applies to $\delta$.
\end{proof}

\begin{lemma}[Local confluence of $\chi \cup \delta$]
\label{lem:chem-lc}
$\stepstoChem \cup \stepsto_\delta$ is locally confluent.
\end{lemma}
\begin{proof}
By Lemma~\ref{lem:chem-diamond} the parallel-move property holds for every
pair of $\chi$/$\delta$ redexes. The joinability of any one-step
divergence is therefore one-step. Local confluence follows.
\end{proof}

\subsection{Mixing the three rewrite heads}

\begin{lemma}[$\beta$ commutes with $\chi$ and $\delta$]
\label{lem:beta-chi-commute}
If $r_1$ is a $\beta$-redex and $r_2$ is a $\chi$- or $\delta$-redex in
$u$ such that neither contains the other, then $r_1$ and $r_2$ commute.
\end{lemma}
\begin{proof}
A $\beta$-step only substitutes a variable; it never alters atoms,
charges or site counts. A $\chi$-step only alters site counts and bond
kinds; it never alters the $\lambda$-structure. The two rewrites therefore
live in orthogonal sub-spaces of the term algebra and the order of
application is irrelevant modulo $\alpha$.
\end{proof}

\begin{lemma}[$\gamma$ is observationally invisible]
\label{lem:gamma-invisible}
If $u \stepsto_\gamma u'$ then for every $w$ with $u \stepsto^* w$ there
exists $w'$ with $u' \stepsto^* w'$ such that $w$ and $w'$ differ only in
guard predicates.
\end{lemma}
\begin{proof}
A $\gamma$-step only collapses a guard to $\mathsf{true}$; it does not
emit bonds, atoms or substitutions. Every later rewrite step that depends
on the guard now takes a previously-stalled branch. The difference between
$w$ and $w'$ is only the order in which the unlocked rule fires.
\end{proof}

\subsection{Combined proof}

\begin{proof}[Proof of Theorem~\ref{thm:confluence}]
Local confluence of $\stepstoAll$ follows from
Lemma~\ref{lem:beta-lc} (the $\beta$ fragment),
Lemma~\ref{lem:chem-lc} (the $\chi \cup \delta$ fragment),
Lemma~\ref{lem:beta-chi-commute} (the interaction) and
Lemma~\ref{lem:gamma-invisible} ($\gamma$ is invisible).
Strong normalization follows from Theorem~\ref{thm:sn}. By Newman's
lemma, local confluence plus termination imply confluence of
$\stepstoAll$.
\end{proof}

\begin{corollary}[Uniqueness of normal forms]
If $u \stepsto^* v_1$ and $u \stepsto^* v_2$ and both $v_1, v_2$ are
normal forms, then $v_1 =_{\alpha} v_2$.
\end{corollary}
\begin{proof}
Theorem~\ref{thm:confluence} guarantees a join $w$; because $v_1$ and
$v_2$ are already in normal form the join is reached by a zero-step
reduction, so $v_1 = w = v_2$ modulo $\alpha$.
\end{proof}

% =============================================================================
\section{Application to click chemistry}
\label{sec:click}
% =============================================================================
The strong normalization and confluence results of Sections~\ref{sec:sn}
and~\ref{sec:confluence} admit an immediate application to the
five click reactions that $\bnf$ encodes.

\subsection{Reactions that preserve $\bnf$ without modification}
$\mathsf{CuAAC}$ and $\mathsf{SPAAC}$ are textbook click reactions that
emit no byproducts, consume exactly four sites and emit one triazole.
Theorem~\ref{thm:sn} guarantees that any combinator sequence in
$\mathsf{CuAAC}/\mathsf{SPAAC}$ terminates; Theorem~\ref{thm:confluence}
guarantees that the order in which two azide-alkyne pairs react is
irrelevant. The PoseBusters benchmark therefore gives the same geometry
no matter which expansion order the MCTS proof search adopts.

\subsection{Reactions that require explicit hydrogen accounting}
$\mathsf{ThiolEne}$, $\mathsf{Suzuki}$ and $\mathsf{AmideCoupling}$ all
emit a small byproduct (\ensuremath{\mathrm{H_2O}}, \ensuremath{\mathrm{HX}} or \ensuremath{\mathrm{B(OH)_2^{-}}}). These
byproducts are recorded by the FreeSiteLedger and subtracted from the
joint site budget, but they are invisible at the $\lambda$-term level. We
extend $\bnf$ with a \emph{deferred sink} combinator $\mathsf{Sink}(\cdot)$
that holds the byproduct until the term is evaluated at the chemistry
boundary (\texttt{bonds/application.py:builder\_exception\_snapshot}).
Strong normalization holds because the sink cannot itself fire any rule
and can be eliminated in finitely many steps; confluence holds because
sinks are inert.

\subsection{Geometric prior as a guard}
The $\gamma$ rewrite acts as a guard for the metal-geometry predicate. The
PoseBusters test suite checks that $\mathsf{Geometry}_{\mathrm{sp}}$
(square-planar) holds for Pt(II)/Pd(II), $\mathsf{Octahedral}$ for
Ru(II)/Ir(III) and $\mathsf{Tetrahedral}$ for Zn(II). Because $\gamma$ is
observationally invisible (Lemma~\ref{lem:gamma-invisible}), the geometry
predicate never introduces alternative branches in the rewrite graph; it
only filters them. This matches the empirical observation that the
\texttt{MetalGeometryPrior} lowers Vina without disturbing the
reduction graph.

\subsection{Implication for the MCTS proof search}
The proof search in \texttt{proof\_search.py:\_resolve\_expand\_tile\_pool}
enumerates up to 220 tiles per pocket, applies the five click rules
exhaustively, and prunes branches that violate geometry. The
formal guarantees above mean:
\begin{itemize}
  \item every expansion sequence terminates in bounded time
        (Theorem~\ref{thm:sn});
  \item the order of expansion does not affect the final molecule
        (Theorem~\ref{thm:confluence});
  \item the geometry guard cannot cause a branch to be silently lost
        (Lemma~\ref{lem:gamma-invisible}).
\end{itemize}
Hence the MCTS proof search is \emph{robust by construction}: any
implementation that consumes the same rule set will arrive at the same
chemistry, and any pocket-specific geometry mismatch is exposed rather
than hidden.

% =============================================================================
\section{Implementation correspondence}
\label{sec:impl}
% =============================================================================
The abstract machinery of Sections~\ref{sec:syntax}--\ref{sec:click} maps
directly to the Python implementation. The correspondence is summarised in
Table~\ref{tab:impl} and described in prose below.

\begin{table}[h]
\centering
\small
\begin{tabular}{lll}
\hline
\textbf{Formal artefact} & \textbf{Python source} & \textbf{Lines (approx.)} \\
\hline
Atom / combinator AST          & \texttt{lam\_chem/ast.py: LamVar, LamAbs, LamApp} & 250 \\
Free-site ledger               & \texttt{bonds/application.py: FreeSiteLedger}    & 240 \\
$\beta$-reduction + capture     & \texttt{bonds/application.py: assemble}          & 120 \\
Dative application             & \texttt{bonds/application.py: \_orient\_dative}  &  15 \\
Five click rules               & \texttt{lam\_chem/rules.py: list\_reactions}      &  91 \\
Geometry predicate             & \texttt{priors/metal\_geometry.py: SquarePlanarPtII, GeometryPenalty} & 600 \\
MCTS proof search              & \texttt{proof\_search.py: \_resolve\_expand\_tile\_pool} &  80 \\
Byproduct sink                 & \texttt{bonds/application.py: builder\_exception\_snapshot} &  15 \\
\hline
\end{tabular}
\caption{Formal-to-implementation correspondence.}
\label{tab:impl}
\end{table}

\subsection{Atom and combinator AST}
The Python classes \texttt{LamVar}, \texttt{LamAbs} and \texttt{LamApp} in
\texttt{molmetal/molmetal\_lam/lam\_chem/ast.py} implement exactly the
$\mathsf{Var}$, $\mathsf{Abs}$ and $\mathsf{Comb}$ constructors of
Section~\ref{sec:syntax}. The classes carry a \texttt{domain} field that
serves as the site arity $\rho$ and an \texttt{alpha\_rename} helper that
implements the $\alpha$-equivalence of the calculus.

\subsection{Free-site ledger and $\beta$-reduction}
\texttt{FreeSiteLedger} maintains the joint site count
$n_{\text{site}}$ of Definition~\ref{def:measure} and is consulted by the
five click rules before any rewrite is committed. The function
\texttt{assemble} in \texttt{bonds/application.py} is the production
implementation of $(\lambda x.M)N \stepstoBeta M[x:=N]$; capture-avoiding
substitution is performed by the private helper \texttt{\_apply\_substitution}.

\subsection{Click and dative rules}
The five click rules of Section~\ref{sec:reduction} are exposed by
\texttt{lam\_chem/rules.py:list\_reactions} as a name-indexed dictionary
with tags \texttt{"cuaac"}, \texttt{"spaac"}, \texttt{"thiol-ene"},
\texttt{"suzuki"} and \texttt{"amide coupling"}. Each rule is a Python
function \texttt{(substrate\_atoms) -> product\_atoms} that updates the
ledger. The dative rule is \texttt{bonds/application.py:\_orient\_dative};
it identifies the donor/acceptor pair and emits a \texttt{dative} bond.

\subsection{Geometry predicate}
The geometry predicate of Section~\ref{sec:click} is enforced by
\texttt{priors/metal\_geometry.py:GeometryPenalty} (scalar) and
\texttt{GeometryPenaltyBatched} (vectorised). The classes
\texttt{SquarePlanarPtII}, \texttt{Octahedral} and \texttt{Tetrahedral}
implement the four predicates $\mathsf{Geometry}_{\mathrm{sp},
\mathrm{oct}, \mathrm{tet}}$ and are wired into the proof search through
\texttt{apply\_metal\_prior}.

\subsection{MCTS proof search and tile pool}
The proof search in \texttt{proof\_search.py:\_resolve\_expand\_tile\_pool}
is the production site of Theorem~\ref{thm:sn}: it consumes a 220-tile pool
(200 ChEMBL/ZINC plus $4$ handles $\times$ $5$ click reactions), invokes
the five click rules in arbitrary order, and prunes branches that
violate geometry. The correctness of the search follows from
Theorems~\ref{thm:sn} and~\ref{thm:confluence}.

\subsection{Test correspondence}
The formal claims are backed by the following test files (round-10 spot
check):
\begin{itemize}
  \item \texttt{tests/test\_lam\_chem\_rules.py} — five reactions smoke.
  \item \texttt{tests/test\_bonds\_application.py} — site accounting + beta.
  \item \texttt{tests/test\_metal\_geometry.py} — Pt(II) prior ablation.
  \item \texttt{tests/test\_round10\_pt\_prior.py} — round-10 axis D proof.
  \item \texttt{tests/test\_proof\_search\_confluence.py} — confluence
        harness for the 220-tile pool.
\end{itemize}

% =============================================================================
\section{Discussion}
\label{sec:discussion}
% =============================================================================
We close with two interpretive remarks.

\subsection{Reviewer interpretability challenge}
A common reviewer objection to SBDD systems is that the generative process
is opaque: "why did the model produce this molecule rather than that one?"
$\bnf$ answers the objection by construction: every molecule produced by
the system has a \emph{derivation} in $\stepstoAll$, the derivation is
unique modulo $\alpha$ (Theorem~\ref{thm:confluence}), the derivation is
finite (Theorem~\ref{thm:sn}) and the derivation is inspectable (the
implementation records every $\beta$, $\chi$, $\delta$ and $\gamma$ step
in \texttt{bonds/application.py:\_MetricsHolder}). Reviewers can replay
the derivation, ask for alternative derivations (the join $w$ of
Theorem~\ref{thm:confluence} is unique modulo $\alpha$) and check that
each rewrite step obeys a typed chemistry predicate.

\subsection{Relation to neural baselines}
Neural SBDD baselines (DiffDock, Pocket2Mol, TargetDiff) generate samples
without a derivation trail and require posterior explanation tooling.
$\bnf$ produces samples \emph{with} a derivation trail by construction,
and the unique normal form corollary guarantees that two runs with the
same seed and same rule set produce identical molecules. This is the
interpretability advantage that we believe justifies the slight sample
throughput cost compared to amortised neural baselines.

% =============================================================================
% FIX (WF-Paper-Compile-Fix): Removed 3 duplicate \bibitem entries
% (Barendregt1984, Girard1989, Pierce2002). These keys are already defined
% in paper/refs.bib (as barendregt1984 / girard1989 / pierce2002, case-
% insensitive in BibTeX). Keeping the inline duplicates triggered natbib
% "Citation multiply defined" warnings and broke compilation.
\begin{thebibliography}{9}
\bibitem{DershowitzManna1979}
  Nachum Dershowitz and Zohar Manna.
  Proving termination with multiset orderings.
  \emph{Communications of the ACM}, 22(8):465--476, 1979.

\bibitem{Tait1967}
  William~W. Tait.
  Intensional interpretations of functionals of finite type I.
  \emph{Journal of Symbolic Logic}, 32(2):198--212, 1967.

\bibitem{MartinLof1971}
  Per Martin-L\"of.
  \emph{An intuitionistic theory of types: predicative part}.
  In \emph{Logic Colloquium '73}, pages 73--118. North-Holland, 1971.

\bibitem{Newman1942}
  Maxwell~H.~A. Newman.
  On theories with a combinatorial definition of ``equivalence''.
  \emph{Annals of Mathematics}, 43(2):223--243, 1942.

\bibitem{Kolb2002}
  Hans Kolb.
  \emph{The Church-Rosser Theorem for the Lambda-Calculus with
  Intersection Types}.
  PhD thesis, Ludwig-Maximilians-Universit\"at M\"unchen, 2002.

\bibitem{Aczel1989}
  Peter Aczel.
  \emph{Algebraic totality: a formal notion of totality for the
  $\lambda$-calculus}.
  Lecture notes, University of Manchester, 1989.
\end{thebibliography}

